\unchapter{Introduction}

\paragraph{Relevance and Problem Context.}
Small and medium-sized e-commerce enterprises (SMEs) operate under acute engineering and financial constraints: compact software engineering teams (typically two to eight developers), constrained monthly cloud infrastructure budgets, and intense commercial pressure to deliver storefront, promotional, and checkout features rapidly \cite{newman2021,ford2022}. Over the past decade, microservices architecture has emerged as the dominant architectural paradigm among hyperscale technology organizations seeking independent service scalability and organizational autonomy \cite{fowler2014,dragoni2017}. However, adopting distributed microservices prematurely introduces substantial operational and runtime overhead—inter-service network serialization, distributed transaction coordination, container orchestration complexity, message broker management, and elevated baseline memory consumption—that can erode an SME's development velocity and infrastructure margin \cite{fowler2015monolith,soldani2018}. Conversely, a poorly structured monolithic application risks tight internal coupling, fragile deployments, and localized performance bottlenecks as catalog volume and concurrent shopper traffic expand \cite{parnas1972,richardson2018}.

\paragraph{Research Object and Subject.}
The \textbf{object} of this study is backend software architecture for small and medium-sized e-commerce web applications. The \textbf{subject} of the study is the empirical trade-off space between a modular monolithic architecture and a distributed microservices architecture across runtime resource footprint, request latency distribution, throughput stability, and functional change propagation.

\paragraph{Aim and Research Tasks.}
The \textbf{aim} of this Bachelor's thesis is to evaluate whether a modular monolithic architecture or a microservices architecture provides a superior engineering and operational balance for SME e-commerce platforms, and to formulate an evidence-based architectural decision framework. To achieve this aim, the following tasks are executed:
\begin{enumerate}[leftmargin=1.5em]
  \item Synthesize theoretical foundations of software modularity, domain-driven design, and distributed systems trade-offs from peer-reviewed literature (Chapter~1).
  \item Analyze SME e-commerce operational constraints, infrastructure cost models, and architectural quality attributes aligned with ISO/IEC~25010:2023 (Chapter~2).
  \item Design and deploy a reproducible containerized experimental testbed comparing two production-grade .NET reference applications: \texttt{eShopOnWeb} (modular monolith) and \texttt{eShopOnContainers} (microservices), and define five standardized functional change scenarios (Chapter~3).
  \item Execute controlled runtime benchmarks (startup readiness, idle and under-load memory, median and tail latency, throughput, and error rates) alongside source-level change propagation analysis across both architectures (Chapter~4).
  \item Triangulate quantitative measurements against published industrial migration case studies to derive a structured SME architectural decision matrix and evaluate threats to validity (Chapter~5).
\end{enumerate}

\paragraph{Research Hypotheses.}
Three testable hypotheses guide the empirical investigation:
\begin{itemize}[leftmargin=1.5em]
  \item \textbf{H1 (Resource Efficiency):} Under SME-scale hardware constraints, the modular monolithic architecture will exhibit at least 15\% lower baseline and under-load memory consumption than the equivalent microservices deployment due to the elimination of per-service runtime processes, API gateways, and message brokers.
  \item \textbf{H2 (Tail Latency and Throughput):} Under concurrent read-heavy storefront workloads, the microservices architecture will achieve lower 95th-percentile ($p_{95}$) tail latency through independent service concurrency and lightweight read models, while maintaining comparable aggregate throughput.
  \item \textbf{H3 (Change Localization vs.\ Cross-Cutting Propagation):} Isolated domain changes will require fewer modified files and deployment units in microservices, whereas cross-cutting domain changes will propagate across significantly more files and deployment units in microservices than in a modular monolith.
\end{itemize}

\paragraph{Methodology and Structure.}
The research employs a mixed-methods empirical software engineering design combining systematic literature review, static source-code dependency and change-impact analysis across concrete C\# solutions, reproducible Docker and \texttt{k6} load benchmarking, and secondary triangulation with empirical industrial migration datasets. Chapter~1 establishes the theoretical framework; Chapter~2 defines SME requirements and cost models; Chapter~3 specifies the experimental testbed and C\# change scenarios; Chapter~4 presents quantitative empirical results; and Chapter~5 synthesizes the decision framework and limitations.
